% Bagian: intuitionistic-logic
% Bab: introduction
% Subbagian: bhk-interpretation

\documentclass[../../../include/open-logic-section]{subfiles}

\begin{document}

\olfileid[id]{int}{int}{bhk}

\olsection{Interpretasi Brouwer--Heyting--Kolmogorov}

\begin{editorial}
  Bukti kesahihan proposisi intuisionistik dengan menggunakan interpretasi
  BHK membingungkan; bukti-bukti tersebut perlu dijelaskan dengan lebih baik.
\end{editorial}

Terdapat suatu interpretasi konstruktif informal bagi konektif intuisionistik,
yang biasanya dikenal sebagai interpretasi Brouwer--Heyting--Kolmogorov.
Interpretasi ini menggunakan gagasan ``konstruksi'', yang dapat dipandang
sebagai bukti konstruktif. (Kita tidak menggunakan istilah ``bukti'' dalam
interpretasi BHK agar tidak tertukar dengan gagasan !!a{derivation} dalam suatu
sistem !!{derivation} formal.) Berdasarkan gagasan intuitif ini, interpretasi
BHK menjelaskan makna konektif intuisionistik.

\begin{enumerate}
\item Kita mengandaikan bahwa kita mengetahui apa yang membentuk suatu
  konstruksi bagi pernyataan atomik.
\item Konstruksi bagi $!A_1 \land !A_2$ ialah pasangan $\tuple{M_1, M_2}$,
  dengan $M_1$ suatu konstruksi bagi~$!A_1$ dan $M_2$ suatu konstruksi
  bagi~$!A_2$.
\item Konstruksi bagi $!A_1 \lor !A_2$ ialah pasangan $\tuple{s, M}$,
  dengan $s$ bernilai~$1$ dan $M$ suatu konstruksi bagi~$!A_1$, atau $s$
  bernilai~$2$ dan $M$ suatu konstruksi bagi~$!A_2$.
\item Konstruksi bagi $!A \lif !B$ ialah fungsi yang mengubah suatu konstruksi
  bagi~$!A$ menjadi konstruksi bagi~$!B$.
\item Tidak ada konstruksi bagi~$\lfalse$ (absurditas).
\item $\lnot !A$ didefinisikan sebagai sinonim bagi $!A \lif \lfalse$.
  Artinya, konstruksi bagi $\lnot !A$ ialah fungsi yang mengubah suatu
  konstruksi bagi~$!A$ menjadi konstruksi bagi~$\lfalse$.
\end{enumerate}

\begin{ex}
Sebagai contoh, ambillah $\lnot \lfalse$. Suatu konstruksi baginya ialah
fungsi yang, ketika menerima sebarang konstruksi bagi $\lfalse$ sebagai
masukan, menghasilkan konstruksi bagi~$\lfalse$ sebagai keluaran. Jelaslah
bahwa fungsi identitas~$\Id{}$ merupakan konstruksi semacam itu: jika diberi
suatu konstruksi~$M$ bagi~$\lfalse$, maka $\Id{}(M) = M$ menghasilkan
konstruksi bagi~$\lfalse$.
\end{ex}

Secara umum, $\lnot !A$ berarti ``Konstruksi bagi $!A$ mustahil.''

\begin{ex}
Mari kita buktikan $!A \lif \lnot \lnot !A$ bagi sebarang proposisi~$!A$,
yakni $!A \lif ((!A \lif \lfalse) \lif \lfalse)$. Konstruksinya harus berupa
fungsi~$f$ yang, ketika diberi konstruksi~$M$ bagi $!A$, mengembalikan
konstruksi~$f(M)$ bagi $(!A \lif \lfalse) \lif \lfalse$. Berikut cara~$f$
membangun konstruksi bagi $(!A \lif \lfalse) \lif \lfalse$: Kita harus
mendefinisikan fungsi $g$ yang, ketika diberi konstruksi~$h$ bagi $!A \lif
\lfalse$ sebagai masukan, menghasilkan konstruksi bagi~$\lfalse$. Kita dapat
mendefinisikan $g$ sebagai berikut: terapkan masukan~$h$ pada konstruksi~$M$
bagi $!A$ (yang telah kita terima sebelumnya). Karena keluaran~$h(M)$ dari $h$
merupakan konstruksi bagi $\lfalse$, maka $f(M)(h) = h(M)$ merupakan konstruksi
bagi~$\lfalse$ jika $M$ merupakan konstruksi bagi~$!A$.
\end{ex}

\begin{ex}
Mari kita berikan konstruksi bagi $\lnot(!A \land \lnot !A)$, yakni $(!A
\land (!A \lif \lfalse)) \lif \lfalse$. Konstruksi ini merupakan fungsi~$f$
yang, ketika diberi konstruksi~$M$ bagi $!A \land (!A \lif \lfalse)$ sebagai
masukan, menghasilkan konstruksi bagi~$\lfalse$. Konstruksi bagi konjungsi
$!B_1 \land !B_2$ ialah pasangan $\tuple{N_1, N_2}$, dengan $N_1$ suatu
konstruksi bagi~$!B_1$ dan $N_2$ suatu konstruksi bagi~$!B_2$. Kita dapat
mendefinisikan fungsi $p_1$ dan $p_2$ yang, dari suatu konstruksi bagi
$!B_1 \land !B_2$, masing-masing mengambil kembali konstruksi bagi $!B_1$
dan $!B_2$:
\begin{align*}
  p_1(\tuple{N_1, N_2}) &= N_1\\
  p_2(\tuple{N_1, N_2}) &= N_2
\end{align*}
Berikut yang dilakukan $f$: Mula-mula, fungsi itu menerapkan $p_1$ pada
masukannya~$M$. Hasilnya ialah suatu konstruksi bagi~$!A$. Kemudian fungsi itu
menerapkan $p_2$ pada~$M$, sehingga menghasilkan konstruksi bagi $!A \lif
\lfalse$. Konstruksi semacam itu sendiri merupakan fungsi~$p_2(M)$ yang, jika
diberi konstruksi bagi~$!A$ sebagai masukan, menghasilkan konstruksi
bagi~$\lfalse$. Dengan kata lain, jika kita menerapkan $p_2(M)$ pada $p_1(M)$,
kita memperoleh konstruksi bagi~$\lfalse$. Jadi, kita dapat mendefinisikan
$f(M) = p_2(M)(p_1(M))$.
\end{ex}

\begin{ex}
Mari kita berikan konstruksi bagi $((!A \land !B) \lif !C) \lif (!A \lif
(!B \lif !C))$, yakni suatu fungsi $f$ yang mengubah konstruksi~$g$ bagi
$(!A \land !B) \lif !C$ menjadi konstruksi bagi $(!A \lif (!B \lif !C))$.
Konstruksi $g$ sendiri merupakan fungsi (dari konstruksi bagi $!A \land !B$
ke konstruksi bagi~$!C$). Keluaran $f(g)$ ialah fungsi~$h_g$ dari konstruksi
bagi $!A$ ke fungsi dari konstruksi bagi~$!B$ ke konstruksi bagi~$!C$.

Baiklah, ini membingungkan. Kita harus membangun fungsi tertentu~$h_g$, yang
akan menjadi keluaran $f$ bagi masukan~$g$. Masukan bagi~$h_g$ ialah
konstruksi~$M$ bagi~$!A$. Keluaran $h_g(M)$ seharusnya berupa fungsi~$k_{g,M}$
dari konstruksi~$N$ bagi~$!B$ ke konstruksi bagi~$!C$. Tetapkan
$k_{g,M}(N) = g(\tuple{M, N})$. Ingat bahwa $\tuple{M, N}$ merupakan
konstruksi bagi~$!A \land !B$. Jadi, $k_{g,M}$ merupakan konstruksi bagi $!B
\lif !C$: fungsi ini memetakan konstruksi~$N$ bagi~$!B$ ke konstruksi
bagi~$!C$. Sekarang tetapkan $h_g(M) = k_{g,M}$. Ini merupakan fungsi yang
memetakan konstruksi~$M$ bagi~$!A$ ke konstruksi $k_{g,M}$ bagi $!B \lif !C$.
Terakhir, tetapkan $f(g) = h_g$. Ini merupakan fungsi yang memetakan
konstruksi $g$ bagi $(!A \land !B) \lif !C$ ke konstruksi bagi $!A \lif (!B
\lif !C)$. Fiuh!{}
\end{ex}

Pernyataan $!A \lor \lnot !A$ disebut Hukum Tiada Jalan Tengah. Kita dapat
membuktikannya bagi beberapa $!A$ tertentu (misalnya, $\lfalse \lor \lnot
\lfalse$), tetapi tidak secara umum. Alasannya ialah bahwa disjungsi
intuisionistik memerlukan konstruksi bagi salah satu disjungnya, sedangkan ada
pernyataan yang saat ini belum dapat dibuktikan ataupun disangkal (misalnya,
konjektur Goldbach). Akan tetapi, hukum tiada jalan tengah juga tidak dapat
Anda sangkal: yakni, $\lnot \lnot (!A \lor \lnot !A)$ berlaku.

\begin{ex}
Untuk membuktikan $\lnot \lnot (!A \lor \lnot !A)$, kita memerlukan fungsi~$f$
yang mengubah konstruksi bagi $\lnot(!A \lor \lnot !A)$, yakni bagi $(!A \lor
(!A \lif \lfalse)) \lif \lfalse$, menjadi konstruksi bagi~$\lfalse$. Dengan
kata lain, kita memerlukan fungsi $f$ sedemikian sehingga $f(g)$ merupakan
konstruksi bagi~$\lfalse$ jika $g$ merupakan konstruksi bagi $\lnot(!A \lor
\lnot !A)$.

Andaikan $g$ merupakan konstruksi bagi $\lnot(!A \lor \lnot !A)$, yakni fungsi
yang mengubah konstruksi bagi $!A \lor \lnot !A$ menjadi konstruksi
bagi~$\lfalse$. Konstruksi bagi $!A \lor \lnot !A$ ialah pasangan
$\tuple{s, M}$, dengan $s = 1$ dan $M$ suatu konstruksi bagi $!A$, atau $s = 2$
dan $M$ suatu konstruksi bagi $\lnot !A$. Misalkan $h_1$ ialah fungsi yang
memetakan konstruksi~$M_1$ bagi $!A$ ke konstruksi bagi $!A \lor \lnot !A$:
fungsi ini memetakan $M_1$ ke $\tuple{1, M_1}$. Misalkan pula $h_2$ ialah
fungsi yang memetakan konstruksi~$M_2$ bagi $\lnot !A$ ke konstruksi bagi
$!A \lor \lnot !A$: fungsi ini memetakan $M_2$ ke $\tuple{2, M_2}$.

Misalkan $k$ adalah $\comp{h_1}{g}$: ini merupakan fungsi yang, jika diberi
konstruksi bagi~$!A$, mengembalikan konstruksi bagi~$\lfalse$; jadi, fungsi ini
merupakan konstruksi bagi~$!A \lif \lfalse$ atau $\lnot !A$. Sekarang misalkan
$l$ adalah $\comp{h_2}{g}$. Ini merupakan fungsi yang, jika diberi konstruksi
bagi $\lnot !A$, menghasilkan konstruksi bagi~$\lfalse$. Karena $k$ merupakan
konstruksi bagi $\lnot !A$, maka $l(k)$ merupakan konstruksi bagi~$\lfalse$.

Secara keseluruhan, yang telah kita lakukan ialah menguraikan cara mengubah
konstruksi~$g$ bagi $\lnot(!A \lor \lnot !A)$ menjadi konstruksi
bagi~$\lfalse$; yakni, fungsi $f$ yang memetakan konstruksi~$g$ bagi
$\lnot(!A \lor \lnot !A)$ ke konstruksi $l(k)$ bagi~$\lfalse$ merupakan
konstruksi bagi $\lnot\lnot(!A \lor \lnot !A)$.
\end{ex}

Sebagaimana dapat Anda lihat, penggunaan interpretasi BHK untuk menunjukkan
kesahihan intuisionistik !!{formula}s dengan cepat menjadi rumit dan
membingungkan. Untunglah, terdapat sistem !!{derivation} yang lebih baik bagi
logika intuisionistik, serta interpretasi semantik yang lebih tepat.
\end{document}
